\documentclass[11pt]{article}

\usepackage[utf8]{inputenc}
\usepackage[T1]{fontenc}
\usepackage{lmodern}
\usepackage{geometry}
\usepackage{amsmath,amssymb,amsfonts,amsthm,mathtools}
\usepackage{bm}
\usepackage{enumitem}
\usepackage{microtype}
\usepackage{hyperref}
\usepackage{titlesec}
\usepackage{booktabs}
\usepackage{array}
\usepackage{setspace}
\usepackage{mathrsfs}
\usepackage{physics}

\geometry{margin=1in}
\hypersetup{colorlinks=true, linkcolor=black, urlcolor=blue, citecolor=black}
\setlist[enumerate]{topsep=0.4em,itemsep=0.35em}
\setlist[itemize]{topsep=0.3em,itemsep=0.25em}
\titleformat{\section}{\large\bfseries}{\thesection.}{0.5em}{}
\titleformat{\subsection}{\normalsize\bfseries}{\thesubsection.}{0.5em}{}
\setstretch{1.06}

\newcommand{\Aether}{\AE{}ther}
\newcommand{\Aflow}{\AE{}ther-flow}
\newcommand{\TheoryName}{\Aflow{} Interpretation of Relativity}
\newcommand{\TheoryProgram}{\Aether{} / \Aflow{} framework}

\newcommand{\CompletedMetric}{g^{\sharp}}
\newcommand{\MatterSpace}{\Psi}

\newcommand{\Car}{\mathrm{car}}
\newcommand{\Cur}{\mathrm{cur}}
\newcommand{\Hom}{\mathrm{hom}}

\newcommand{\ForSpace}[1]{\mathcal I_{#1}^{\mathrm{for}}}
\newcommand{\PrimShell}{\mathcal S_{\mathrm{prim}}}
\newcommand{\MetricOut}{\mathscr G_{\mathrm{out}}}
\newcommand{\MatterOut}{\mathscr T_{\mathrm{out}}}

\newcommand{\ProtoSectionPackage}{\mathfrak Q^{\mathrm{sub,proto,secgen}}}
\newcommand{\ProtoSectionBody}[1]{\mathcal B_{#1}^{\mathrm{psec}}}
\newcommand{\ProtoSectionShell}[1]{\mathcal Z_{#1}^{\mathrm{psec}}}

\newcommand{\PreProtoSectionPackage}{\mathfrak R^{\mathrm{sub,proto,presec}}}
\newcommand{\PreProtoSectionMinSectionCore}{\mathfrak R^{\mathrm{sub,proto,presec,sec}}}
\newcommand{\PreProtoSectionResidualPackage}{\mathfrak R^{\mathrm{sub,proto,presec,res}}}

\newcommand{\PrePreProtoSectionPullbackCore}{\mathfrak S^{\mathrm{sub,proto,pppresec,pb}}}
\newcommand{\PrePreProtoSectionVertical}[1]{V_{#1}^{\mathrm{pppresec}}}
\newcommand{\PrePreProtoSectionResidualBody}[1]{\mathcal D_{#1}^{\mathrm{pppresec}}}
\newcommand{\PrePreProtoSectionBodyResponse}[1]{\mathscr R_{#1}^{\mathrm{pppresec,pb,body}}}

\newcommand{\PullbackOperatorCore}{\mathfrak S^{\mathrm{sub,proto,pppresec,pb,op}}}
\newcommand{\PullbackBodyOperatorCore}{\mathfrak S^{\mathrm{sub,proto,pppresec,pb,bodyop}}}
\newcommand{\PullbackBodyResponseCore}{\mathfrak S^{\mathrm{sub,proto,pppresec,pb,bodyresp}}}

\newcommand{\ProjectionOperatorCore}{\mathfrak S^{\mathrm{disp,resp,projop}}}
\newcommand{\SubstrateCarrier}[1]{U_{#1}^{\mathrm{disp,resp}}}
\newcommand{\SubstrateBody}[1]{\mathcal U_{#1}^{\mathrm{disp,resp}}}
\newcommand{\SubstrateProjection}[1]{P_{#1}^{\mathrm{disp,resp}}}
\newcommand{\SubstrateOperator}[1]{K_{#1}^{\mathrm{disp,resp}}}
\newcommand{\SubstrateLift}[1]{J_{#1}^{\mathrm{disp,resp}}}
\newcommand{\CanonicalLift}[1]{J_{#1}^{\mathrm{disp,resp,can}}}
\newcommand{\SubstrateKernel}[1]{N_{#1}^{\mathrm{disp,resp}}}
\newcommand{\SubstrateEnergy}[1]{\mathscr E_{#1}^{\mathrm{disp,resp}}}
\newcommand{\CoreForcedBodyResponse}[1]{\widehat{\mathscr R}_{#1}^{\mathrm{disp,resp,projop}}}
\newcommand{\SchurResponse}[1]{M_{#1}^{\mathrm{disp,resp}}}

\newtheorem{definition}{Definition}
\newtheorem{proposition}{Proposition}
\newtheorem{remark}{Remark}
\newtheorem{corollary}{Corollary}
\newtheorem{theorem}{Theorem}

\title{\textbf{The \TheoryName{}}\\[0.4em]
\large Primitive-Reservoir Observer-Reduction\\
\large Section-Centered Hidden-Response\\
\large Displacement-Response Substrate Law\\
\large Collapse to a Projection-Operator Core}
\author{}
\date{}

\begin{document}

\maketitle

\begin{abstract}
The current active-primary theorem already removed the observer-relevant pullback body-response core \(\PullbackBodyResponseCore\) as the live unexplained stopping point on the recorded section-centered hidden-response route by deriving that core from a benchmark-faithful current-body-realizing section-centered displacement-response substrate law. The honest next move below that frontier is therefore no longer another same-law relay. It is to derive that law itself from still more primitive explicit substrate structure or, if that deeper forcing package is not yet cleanly in hand, to prove a genuinely smaller theorem inside the law.

The present note takes that smaller route. It proves that the full current displacement-response substrate law still carries benchmark-facing presentational redundancy: once the current-body-realizing substrate body, the surjective displacement projection, the coercive positive self-adjoint substrate-response operator, the exact-shell discipline, and benchmark faithfulness are fixed, the \(K\)-orthogonal minimal-lift splitting is no longer separate primitive data. It is uniquely forced by an explicit weighted left-inverse formula
\[
\CanonicalLift{\sigma}(y)
=
\SubstrateOperator{\sigma}(y)^{-1}\SubstrateProjection{\sigma}(y)^{*}
\Bigl(
\SubstrateProjection{\sigma}(y)\SubstrateOperator{\sigma}(y)^{-1}\SubstrateProjection{\sigma}(y)^{*}
\Bigr)^{-1}.
\]
Equivalently, the full displacement-response substrate law collapses benchmark-facingly to a smaller current-body-realizing projection-operator core \(\ProjectionOperatorCore\).

The benchmark boundary remains fixed throughout. The one operative metric is still exactly \(\CompletedMetric\), the ordinary matter route is unchanged, there is no second operative metric, the low-energy matter law is not enlarged, and no stronger promotion language is licensed. The positive result is therefore a genuine smaller theorem inside the current substrate-law frontier. What remains open is to derive or force that smaller projection-operator core itself, or some nontrivial constituent or compatibility relation inside it, from still more primitive explicit substrate structure.
\end{abstract}

\section{Introduction}

The active derivational program of \TheoryName{} has again reached the point where depth alone is not the honest next move. The current section-centered hidden-response frontier theorem already showed that the observer-relevant pullback body-response core \(\PullbackBodyResponseCore\) is not primitive: it is forced by a benchmark-faithful current-body-realizing section-centered displacement-response substrate law. Once that theorem is in hand, another manuscript that merely carries along the same substrate-law data and reproduces the same current downstream benchmark one-metric package no longer removes the live burden.

What remains honest is either to derive that current substrate law from still more primitive explicit substrate structure or to prove a strictly smaller theorem inside the law. The present note takes that internal route. Its claim is that the separate minimal-lift family inside the current displacement-response substrate law is not additional benchmark-facing freedom. It is already forced by the displacement projection and the coercive substrate-response operator once the current-body-realizing substrate body and exact-shell discipline are fixed.

\paragraph{Framework and claim boundary.}
\TheoryProgram{} denotes the broader \Aether{} / \Aflow{} framework built on the claims that the \Aether{} is the underlying four-dimensional substrate of reality, the \Aflow{} is its intrinsic ordered motion, observed three-dimensional space is the local experiential slice of that deeper substrate, S-time is the experienced order of change arising from matter, light, and the \Aflow{}, and observed expansion is the three-dimensional appearance of deeper four-dimensional ordered motion. Within that framework, \TheoryName{} denotes the exact-closure theory in which the effective gravitational dynamics are adopted to be exactly Einsteinian with universal matter coupling. In the active sequence, \emph{adoption} means use of that established relativistic dynamics without claiming substrate derivation, while \emph{derivation} is reserved for a first-principles recovery from explicit substrate variables.

The present note stays strictly inside that boundary. It does not introduce a second operative metric, it does not enlarge the low-energy matter law, and it does not claim a completed first-principles derivation of GR from final substrate microphysics. Its narrower benchmark-facing role is to identify a smaller live bridge object strictly inside the current section-centered displacement-response substrate-law frontier.

\section{Fixed Inputs from the Current Frontier}

\begin{definition}[Recorded primitive continuation data]
\label{def:fixed}
Fix the following continuation data from the active primitive-reservoir observer-reduction line.
\begin{enumerate}
    \item For each sector label \(\sigma\in\{\Car,\Cur,\Hom\}\), a primitive forcing shell
    \[
    \ForSpace{\sigma}.
    \]
    \item The product shell
    \[
    \PrimShell
    \equiv
    \ForSpace{\Car}\times \ForSpace{\Cur}\times \ForSpace{\Hom}.
    \]
    \item The recorded metric and ordinary-matter outputs
    \[
    \MetricOut:\PrimShell\to\{\text{observer metrics}\},
    \qquad
    \MatterOut:\PrimShell\times\MatterSpace\to\{\text{observer stress tensors}\},
    \]
    satisfying
    \begin{equation}
    \MetricOut(\mathbf w)=\CompletedMetric,
    \qquad
    \MatterOut(\mathbf w,\psi)
    =
    T_{\mu\nu}^{\mathrm{matter}}[\CompletedMetric,\psi]
    \label{eq:fixedoutputs}
    \end{equation}
    for every \(\mathbf w\in\PrimShell\) and every matter configuration \(\psi\in\MatterSpace\).
\end{enumerate}
\end{definition}

\begin{definition}[Recorded benchmark base and downstream chain]
\label{def:downstream}
Fix \(\sigma\in\{\Car,\Cur,\Hom\}\). The active line already fixes the following data and consequences.
\begin{enumerate}
    \item the current proto-section benchmark base \(\ProtoSectionPackage\), in particular the recorded proto-section body \(\ProtoSectionBody{\sigma}\) and shell \(\ProtoSectionShell{\sigma}\);
    \item the current observer-relevant pullback body-response core \(\PullbackBodyResponseCore\), the current observer-relevant pullback body-operator core \(\PullbackBodyOperatorCore\), the current observer-relevant pullback-operator core \(\PullbackOperatorCore\), the current observer-relevant pullback-quadratic core \(\PrePreProtoSectionPullbackCore\), the current section-centered residual-normal-form realization \(\PreProtoSectionResidualPackage\), the current observer-relevant pre-proto-section minimizer-section core \(\PreProtoSectionMinSectionCore\), and the current benchmark-faithful pre-proto-section package \(\PreProtoSectionPackage\);
    \item the previously recorded fact that, once a benchmark-faithful observer-relevant pullback body-response core over the fixed proto-section benchmark base is available, the current body-operator-core realization, the current pullback-operator-core realization, the current pullback-quadratic-core realization, the current section-centered residual-normal-form realization, the current pre-proto-section package, the current minimizer-section core, and the previously recorded continuity-Hessian compressed core and benchmark one-metric observer-package consequences follow unchanged;
    \item benchmark faithfulness, meaning preservation of the benchmark outputs \eqref{eq:fixedoutputs}, no second operative metric, no enlarged low-energy matter law, no change to the ordinary matter route, and no premature promotion from adoption to completed derivation.
\end{enumerate}
\end{definition}

\begin{definition}[Current observer-relevant pullback body-response core]
\label{def:bodyresp}
Fix \(\sigma\in\{\Car,\Cur,\Hom\}\). A \emph{benchmark-faithful observer-relevant pullback body-response core} \(\PullbackBodyResponseCore\) for sector \(\sigma\) consists of:
\begin{enumerate}
    \item a compact convex section-centered displacement body
    \[
    \PrePreProtoSectionResidualBody{\sigma}
    \subset
    \ProtoSectionBody{\sigma}\times\PrePreProtoSectionVertical{\sigma}
    \]
    whose fiber
    \[
    \PrePreProtoSectionResidualBody{\sigma}(y)
    \equiv
    \left\{
    v\in\PrePreProtoSectionVertical{\sigma}:(y,v)\in\PrePreProtoSectionResidualBody{\sigma}
    \right\}
    \]
    is nonempty, compact, convex, contains \(0\), and has nonempty interior for every \(y\in\ProtoSectionBody{\sigma}\);
    \item the body-restricted pullback response
    \[
    \PrePreProtoSectionBodyResponse{\sigma}:
    \PrePreProtoSectionResidualBody{\sigma}\to[0,\infty),
    \]
    smooth in \(y\), fiberwise quadratic in \(v\), and satisfying
    \[
    \PrePreProtoSectionBodyResponse{\sigma}(y,v)=0
    \qquad\Longleftrightarrow\qquad
    v=0
    \]
    on every fiber;
    \item benchmark faithfulness in the sense of Definition \ref{def:downstream}(4).
\end{enumerate}
\end{definition}

\begin{proposition}[What the current frontier already fixes]
\label{prop:frontier}
On the recorded section-centered hidden-response route, the current frontier already fixes two benchmark-facing facts.
\begin{enumerate}
    \item The observer-relevant pullback body-response core \(\PullbackBodyResponseCore\) is the unique minimal benchmark-facing datum at current scope up to body-response-faithful equivalence.
    \item All current benchmark-facing downstream uses already factor through the current body-response core and the fixed proto-section benchmark base.
\end{enumerate}
\end{proposition}

\begin{proof}
Item (1) is exactly the current body-response-core necessity / uniqueness theorem. Item (2) is exactly the current theorem deriving that same core from the section-centered displacement-response substrate law together with the earlier collapse theorems that already routed all present downstream benchmark-facing uses through the smaller body-response core.
\end{proof}

\section{The Current Substrate-Law Frontier}

\begin{definition}[Current-body-realizing section-centered displacement-response substrate law]
\label{def:law}
Fix \(\sigma\in\{\Car,\Cur,\Hom\}\). A \emph{benchmark-faithful current-body-realizing section-centered displacement-response substrate law} for sector \(\sigma\) consists of the following data.
\begin{enumerate}
    \item A finite-dimensional Euclidean substrate carrier
    \[
    \SubstrateCarrier{\sigma}
    \]
    and a compact convex section-centered substrate body
    \[
    \SubstrateBody{\sigma}
    \subset
    \ProtoSectionBody{\sigma}\times\SubstrateCarrier{\sigma}
    \]
    whose fiber
    \[
    \SubstrateBody{\sigma}(y)
    \equiv
    \left\{
    u\in\SubstrateCarrier{\sigma}:(y,u)\in\SubstrateBody{\sigma}
    \right\}
    \]
    is nonempty, compact, convex, contains \(0\), and has nonempty interior for every \(y\in\ProtoSectionBody{\sigma}\).
    \item A smooth family of surjective linear displacement projections
    \[
    \SubstrateProjection{\sigma}(y):
    \SubstrateCarrier{\sigma}\to\PrePreProtoSectionVertical{\sigma},
    \qquad
    y\in\ProtoSectionBody{\sigma}.
    \]
    \item A smooth family of positive self-adjoint substrate-response operators
    \[
    \SubstrateOperator{\sigma}(y):
    \SubstrateCarrier{\sigma}\to\SubstrateCarrier{\sigma},
    \qquad
    y\in\ProtoSectionBody{\sigma},
    \]
    satisfying the uniform coercive bound
    \begin{equation}
    \ev{u,\SubstrateOperator{\sigma}(y)u}
    \ge
    \mu_{\sigma}\,\|u\|^{2}
    \qquad
    \text{for all }
    y\in\ProtoSectionBody{\sigma},
    \ \ u\in\SubstrateCarrier{\sigma}
    \label{eq:coercive}
    \end{equation}
    for some constant \(\mu_{\sigma}>0\).
    \item A smooth family of linear minimal lifts
    \[
    \SubstrateLift{\sigma}(y):
    \PrePreProtoSectionVertical{\sigma}\to\SubstrateCarrier{\sigma}
    \]
    satisfying, for every \(y\in\ProtoSectionBody{\sigma}\),
    \begin{enumerate}
        \item
        \[
        \SubstrateProjection{\sigma}(y)\circ\SubstrateLift{\sigma}(y)
        =
        \mathrm{id}_{\PrePreProtoSectionVertical{\sigma}};
        \]
        \item every \(u\in\SubstrateCarrier{\sigma}\) admits the unique decomposition
        \[
        u
        =
        \SubstrateLift{\sigma}(y)\bigl(\SubstrateProjection{\sigma}(y)u\bigr)
        +
        k,
        \qquad
        k\in\SubstrateKernel{\sigma}(y)\equiv\ker\SubstrateProjection{\sigma}(y);
        \]
        \item the lift is \(\SubstrateOperator{\sigma}(y)\)-orthogonal to the kernel, namely
        \[
        \ev{
        \SubstrateLift{\sigma}(y)v,
        \SubstrateOperator{\sigma}(y)k
        }
        =
        0
        \]
        for every \(v\in\PrePreProtoSectionVertical{\sigma}\) and every \(k\in\SubstrateKernel{\sigma}(y)\).
    \end{enumerate}
    \item Exact-shell discipline, meaning that the projected zero slice of the substrate body records exactly the zero section over the fixed proto-section shell:
    \[
    \left\{
    y\in\ProtoSectionBody{\sigma}:
    \exists u\in\SubstrateBody{\sigma}(y)
    \text{ with }
    \SubstrateProjection{\sigma}(y)u=0
    \right\}
    =
    \ProtoSectionShell{\sigma}.
    \]
    \item Current-body realization:
    \begin{equation}
    \PrePreProtoSectionResidualBody{\sigma}(y)
    =
    \SubstrateProjection{\sigma}(y)\bigl(\SubstrateBody{\sigma}(y)\bigr)
    \qquad
    \text{for every }
    y\in\ProtoSectionBody{\sigma}.
    \label{eq:realizing}
    \end{equation}
    \item Benchmark faithfulness in the sense of Definition \ref{def:downstream}(4).
\end{enumerate}
\end{definition}

\begin{remark}
Definition \ref{def:law} is exactly the current frontier package. The present note asks whether the separately listed minimal-lift family \(\SubstrateLift{\sigma}\) still carries benchmark-facing freedom once the rest of that law is fixed.
\end{remark}

\section{Projection-Operator Core}

\begin{definition}[Current-body-realizing section-centered displacement-response projection-operator core]
\label{def:core}
Fix \(\sigma\in\{\Car,\Cur,\Hom\}\). A \emph{benchmark-faithful current-body-realizing section-centered displacement-response projection-operator core} \(\ProjectionOperatorCore\) for sector \(\sigma\) consists of the data of Definition \ref{def:law}(1), (2), (3), (5), (6), and (7), but \emph{not} the separately listed minimal-lift family of Definition \ref{def:law}(4).

Equivalently, \(\ProjectionOperatorCore\) retains:
\begin{enumerate}
    \item the finite-dimensional Euclidean substrate carrier \(\SubstrateCarrier{\sigma}\) and compact convex section-centered substrate body \(\SubstrateBody{\sigma}\);
    \item the surjective displacement projection \(\SubstrateProjection{\sigma}(y)\);
    \item the coercive positive self-adjoint substrate-response operator \(\SubstrateOperator{\sigma}(y)\);
    \item the exact-shell discipline;
    \item the current-body realization \eqref{eq:realizing};
    \item benchmark faithfulness.
\end{enumerate}
\end{definition}

\begin{remark}
Definition \ref{def:core} is strictly smaller than Definition \ref{def:law}. It asks whether the benchmark-facing content of the current substrate-law frontier already lies in the substrate body, the displacement projection, the coercive response operator, and their exact-shell and current-body compatibility, with the minimal-lift family treated as derived rather than separately primitive.
\end{remark}

\section{Canonical Minimal Lift from the Projection-Operator Core}

\begin{proposition}[Positive Schur-response operator]
\label{prop:schur}
Fix a current-body-realizing projection-operator core in the sense of Definition \ref{def:core}. For each \(y\in\ProtoSectionBody{\sigma}\), define
\[
\SchurResponse{\sigma}(y)
\equiv
\SubstrateProjection{\sigma}(y)
\SubstrateOperator{\sigma}(y)^{-1}
\SubstrateProjection{\sigma}(y)^*:
\PrePreProtoSectionVertical{\sigma}\to\PrePreProtoSectionVertical{\sigma},
\]
where \((\cdot)^*\) denotes Euclidean adjoint with respect to the fixed inner products on \(\SubstrateCarrier{\sigma}\) and \(\PrePreProtoSectionVertical{\sigma}\). Then \(\SchurResponse{\sigma}(y)\) is positive self-adjoint and invertible for every \(y\in\ProtoSectionBody{\sigma}\).
\end{proposition}

\begin{proof}
Self-adjointness is immediate from the self-adjointness of \(\SubstrateOperator{\sigma}(y)^{-1}\). For any \(0\neq w\in\PrePreProtoSectionVertical{\sigma}\),
\[
\ev{w,\SchurResponse{\sigma}(y)w}
=
\ev{
\SubstrateOperator{\sigma}(y)^{-1/2}\SubstrateProjection{\sigma}(y)^*w,
\SubstrateOperator{\sigma}(y)^{-1/2}\SubstrateProjection{\sigma}(y)^*w
}
\ge 0.
\]
If this quantity vanishes, then \(\SubstrateProjection{\sigma}(y)^*w=0\). Because \(\SubstrateProjection{\sigma}(y)\) is surjective, its Euclidean adjoint is injective, so \(w=0\), a contradiction. Thus \(\SchurResponse{\sigma}(y)\) is positive definite and therefore invertible on the finite-dimensional space \(\PrePreProtoSectionVertical{\sigma}\).
\end{proof}

\begin{definition}[Canonical lift and core-forced response]
\label{def:canonical}
Fix a current-body-realizing projection-operator core. For each \(y\in\ProtoSectionBody{\sigma}\), define the \emph{canonical minimal lift}
\begin{equation}
\CanonicalLift{\sigma}(y)
\equiv
\SubstrateOperator{\sigma}(y)^{-1}\SubstrateProjection{\sigma}(y)^*
\SchurResponse{\sigma}(y)^{-1}
:
\PrePreProtoSectionVertical{\sigma}\to\SubstrateCarrier{\sigma}.
\label{eq:canonical}
\end{equation}
Define the quadratic substrate-response energy by
\[
\SubstrateEnergy{\sigma}(y,u)
\equiv
\frac{1}{2}
\ev{u,\SubstrateOperator{\sigma}(y)u}
\]
and the \emph{core-forced body response} on the current displacement body by
\begin{equation}
\CoreForcedBodyResponse{\sigma}(y,v)
\equiv
\frac{1}{2}
\ev{
\CanonicalLift{\sigma}(y)v,
\SubstrateOperator{\sigma}(y)\CanonicalLift{\sigma}(y)v
}
\qquad
\text{for }
(y,v)\in\PrePreProtoSectionResidualBody{\sigma}.
\label{eq:coreresponse}
\end{equation}
\end{definition}

\begin{proposition}[Canonical lift properties]
\label{prop:canonical}
Fix a current-body-realizing projection-operator core in the sense of Definition \ref{def:core}. Then the canonical lift of Definition \ref{def:canonical} satisfies, for every \(y\in\ProtoSectionBody{\sigma}\),
\begin{enumerate}
    \item
    \[
    \SubstrateProjection{\sigma}(y)\circ\CanonicalLift{\sigma}(y)
    =
    \mathrm{id}_{\PrePreProtoSectionVertical{\sigma}};
    \]
    \item
    \[
    \ev{
    \CanonicalLift{\sigma}(y)v,
    \SubstrateOperator{\sigma}(y)k
    }
    =
    0
    \]
    for every \(v\in\PrePreProtoSectionVertical{\sigma}\) and every \(k\in\SubstrateKernel{\sigma}(y)\);
    \item every \(u\in\SubstrateCarrier{\sigma}\) admits the unique decomposition
    \[
    u
    =
    \CanonicalLift{\sigma}(y)\bigl(\SubstrateProjection{\sigma}(y)u\bigr)
    +
    k,
    \qquad
    k\in\SubstrateKernel{\sigma}(y);
    \]
    \item \(\CanonicalLift{\sigma}(y)\) is linear in the displacement argument and smooth in \(y\).
\end{enumerate}
\end{proposition}

\begin{proof}
Item (1) follows directly from Definition \ref{def:canonical}:
\[
\SubstrateProjection{\sigma}(y)\CanonicalLift{\sigma}(y)
=
\SchurResponse{\sigma}(y)\SchurResponse{\sigma}(y)^{-1}
=
\mathrm{id}.
\]
For item (2), let \(k\in\SubstrateKernel{\sigma}(y)\). Then
\[
\ev{
\CanonicalLift{\sigma}(y)v,
\SubstrateOperator{\sigma}(y)k
}
=
\ev{
\SubstrateProjection{\sigma}(y)^*
\SchurResponse{\sigma}(y)^{-1}v,
k
}
=
\ev{
\SchurResponse{\sigma}(y)^{-1}v,
\SubstrateProjection{\sigma}(y)k
}
=
0.
\]
For item (3), if \(u\in\SubstrateCarrier{\sigma}\), set
\[
k
\equiv
u-\CanonicalLift{\sigma}(y)\bigl(\SubstrateProjection{\sigma}(y)u\bigr).
\]
Applying \(\SubstrateProjection{\sigma}(y)\) and using item (1) gives \(\SubstrateProjection{\sigma}(y)k=0\), so \(k\in\SubstrateKernel{\sigma}(y)\). Uniqueness follows because the right inverse of item (1) fixes the non-kernel part.

Linearity in the displacement argument is immediate from \eqref{eq:canonical}. Smoothness in \(y\) follows because \(\SubstrateProjection{\sigma}(y)\), \(\SubstrateOperator{\sigma}(y)\), their adjoints and inverses, and \(\SchurResponse{\sigma}(y)^{-1}\) all vary smoothly with \(y\).
\end{proof}

\begin{proposition}[Unique energy-minimizing lift]
\label{prop:min}
Fix a current-body-realizing projection-operator core in the sense of Definition \ref{def:core}. For every \(y\in\ProtoSectionBody{\sigma}\) and every \(v\in\PrePreProtoSectionVertical{\sigma}\), the canonical lift \(\CanonicalLift{\sigma}(y)v\) is the unique minimizer of the constrained quadratic problem
\begin{equation}
\min\left\{
\SubstrateEnergy{\sigma}(y,u):
u\in\SubstrateCarrier{\sigma},
\ \SubstrateProjection{\sigma}(y)u=v
\right\}.
\label{eq:minproblem}
\end{equation}
Consequently, for every \(u\in\SubstrateCarrier{\sigma}\) with \(\SubstrateProjection{\sigma}(y)u=v\),
\begin{equation}
\begin{aligned}
\ev{u,\SubstrateOperator{\sigma}(y)u}
&=
\ev{
\CanonicalLift{\sigma}(y)v,
\SubstrateOperator{\sigma}(y)\CanonicalLift{\sigma}(y)v
}
+
\ev{k,\SubstrateOperator{\sigma}(y)k},\\
k
&\equiv
u-\CanonicalLift{\sigma}(y)v
\in
\SubstrateKernel{\sigma}(y),
\end{aligned}
\label{eq:energysplit}
\end{equation}
and the second term vanishes if and only if \(k=0\).
\end{proposition}

\begin{proof}
Let \(u\in\SubstrateCarrier{\sigma}\) satisfy \(\SubstrateProjection{\sigma}(y)u=v\). By Proposition \ref{prop:canonical}(3),
\[
u
=
\CanonicalLift{\sigma}(y)v+k,
\qquad
k\in\SubstrateKernel{\sigma}(y).
\]
Using Proposition \ref{prop:canonical}(2),
\[
\begin{aligned}
\ev{u,\SubstrateOperator{\sigma}(y)u}
&=
\ev{
\CanonicalLift{\sigma}(y)v,
\SubstrateOperator{\sigma}(y)\CanonicalLift{\sigma}(y)v
}\\
&\quad+
2\ev{
\CanonicalLift{\sigma}(y)v,
\SubstrateOperator{\sigma}(y)k
}
+
\ev{k,\SubstrateOperator{\sigma}(y)k}.
\end{aligned}
\]
reduces to \eqref{eq:energysplit}. By the coercive bound \eqref{eq:coercive}, the last term is nonnegative and vanishes only when \(k=0\). Hence the unique constrained minimizer is \(u=\CanonicalLift{\sigma}(y)v\).
\end{proof}

\section{Collapse of the Current Substrate-Law Frontier}

\begin{theorem}[Displacement-response substrate law collapse to a projection-operator core]
\label{thm:main}
Fix the primitive continuation data of Definition \ref{def:fixed}, the recorded benchmark base and downstream chain of Definition \ref{def:downstream}, and for each sector \(\sigma\in\{\Car,\Cur,\Hom\}\) a benchmark-faithful current-body-realizing section-centered displacement-response substrate law in the sense of Definition \ref{def:law}. Then:
\begin{enumerate}
    \item forgetting the separately listed minimal-lift family \(\SubstrateLift{\sigma}\) leaves a benchmark-faithful current-body-realizing projection-operator core \(\ProjectionOperatorCore\) in the sense of Definition \ref{def:core};
    \item the original minimal-lift family is uniquely forced by that smaller core and equals the canonical lift of Definition \ref{def:canonical},
    \[
    \SubstrateLift{\sigma}(y)=\CanonicalLift{\sigma}(y)
    \qquad
    \text{for every }
    y\in\ProtoSectionBody{\sigma};
    \]
    \item the body-restricted pullback response forced by the full substrate law is therefore already forced by the smaller projection-operator core,
    \[
    \CoreForcedBodyResponse{\sigma}(y,v)
    =
    \frac{1}{2}
    \ev{
    \CanonicalLift{\sigma}(y)v,
    \SubstrateOperator{\sigma}(y)\CanonicalLift{\sigma}(y)v
    };
    \]
    \item every benchmark-faithful full substrate law compatible with the same smaller projection-operator core yields the same canonical lift, the same body-response core, and hence the same downstream pullback body-operator core, pullback-operator core, pullback-quadratic core, section-centered residual-normal-form realization, pre-proto-section package, continuity-Hessian compressed core, and benchmark one-metric observer package.
\end{enumerate}
\end{theorem}

\begin{proof}
Item (1) is immediate from Definition \ref{def:law}: the smaller core simply forgets the separately listed data of item (4) while keeping items (1), (2), (3), (5), (6), and (7).

For item (2), fix \(y\in\ProtoSectionBody{\sigma}\) and \(v\in\PrePreProtoSectionVertical{\sigma}\). By Definition \ref{def:law}(4), the original lift \(\SubstrateLift{\sigma}(y)v\) lies in the affine fiber
\[
\left\{
u\in\SubstrateCarrier{\sigma}:
\SubstrateProjection{\sigma}(y)u=v
\right\}
\]
and is \(\SubstrateOperator{\sigma}(y)\)-orthogonal to the kernel \(\SubstrateKernel{\sigma}(y)\). Repeating the energy decomposition used in Proposition \ref{prop:min} shows that \(\SubstrateLift{\sigma}(y)v\) is the unique minimizer of \eqref{eq:minproblem}. By Proposition \ref{prop:min}, that unique minimizer is \(\CanonicalLift{\sigma}(y)v\). Hence \(\SubstrateLift{\sigma}(y)=\CanonicalLift{\sigma}(y)\).

Item (3) is just Definition \ref{def:canonical} together with item (2): the response induced by the full law equals the response induced by the canonical lift forced by the smaller core.

For item (4), any compatible full law has the same data retained in Definition \ref{def:core}. Item (2) therefore forces the same lift, and item (3) forces the same body-response core on the current displacement body. Proposition \ref{prop:frontier}(2) then gives the same downstream benchmark-facing consequences.
\end{proof}

\begin{corollary}[The minimal-lift splitting is no longer primitive]
\label{cor:nonprimitive}
Under the hypotheses of Theorem \ref{thm:main}, the separate minimal-lift family inside the current section-centered displacement-response substrate law is no longer a primitive stopping point on the recorded route. It is uniquely forced by the smaller current-body-realizing projection-operator core \(\ProjectionOperatorCore\).
\end{corollary}

\begin{proof}
This is exactly item (2) of Theorem \ref{thm:main}.
\end{proof}

\begin{corollary}[The current body-response core already follows from the smaller core]
\label{cor:bodyresp}
Under the hypotheses of Theorem \ref{thm:main}, the pair
\[
\left(
\PrePreProtoSectionResidualBody{\sigma},
\CoreForcedBodyResponse{\sigma}
\right)
\]
is a benchmark-faithful observer-relevant pullback body-response core over the fixed proto-section benchmark base, and it is body-response-faithfully equivalent to the unique minimal benchmark-facing body-response datum on the recorded route.
\end{corollary}

\begin{proof}
By Definition \ref{def:core}, the smaller core retains the current-body realization \eqref{eq:realizing}. Proposition \ref{prop:min} shows that \(\CoreForcedBodyResponse{\sigma}\) is a smooth fiberwise quadratic form on the current displacement body that vanishes exactly on the zero displacement. Thus the pair is a benchmark-faithful body-response core. Proposition \ref{prop:frontier}(1) then gives the stated body-response-faithful equivalence.
\end{proof}

\begin{corollary}[Primary post-collapse burden]
\label{cor:next}
After Theorem \ref{thm:main}, the next primary manuscript below the current frontier must do at least one of the following:
\begin{enumerate}
    \item derive or force the current-body-realizing projection-operator core \(\ProjectionOperatorCore\) itself from still more primitive explicit substrate structure;
    \item derive or force a nontrivial constituent of that smaller core, such as the substrate body, the displacement projection, the coercive substrate-response operator, or the exact-shell/current-body compatibility relation;
    \item prove a sharper necessity, uniqueness, or compatibility theorem for a nontrivial constituent or relation inside that smaller core;
    \item do one of the previous three while preserving the same benchmark one-metric package, the same ordinary matter route, and the same claim boundary.
\end{enumerate}
Another manuscript that merely reinstates the canonical minimal-lift splitting as if it were still separate primitive data, merely replays the full displacement-response substrate law, or re-expands back to the larger pullback-body-operator, pullback-operator, pullback-quadratic, hidden-response, residual-normal-form, or pre-proto-section presentations is not the honest default next move.
\end{corollary}

\begin{proof}
By Corollary \ref{cor:nonprimitive}, the separate minimal-lift family is no longer the live unexplained stopping point. The live burden has therefore moved to the smaller projection-operator core that now forces it, or to a strictly smaller theorem inside that core.
\end{proof}

\begin{proposition}[The benchmark package remains fixed]
\label{prop:benchmark}
The present collapse theorem preserves the benchmark one-metric package \eqref{eq:fixedoutputs}. In particular, the operative metric remains exactly \(\CompletedMetric\), the ordinary matter route is unchanged, there is no second operative metric, and the low-energy matter law is not enlarged.
\end{proposition}

\begin{proof}
This is part of benchmark faithfulness in Definitions \ref{def:downstream}(4) and \ref{def:core}. The present note changes only the upstream status of the current section-centered displacement-response substrate-law frontier.
\end{proof}

\section{Conclusion}

The positive result of this note is a genuine smaller theorem inside the current section-centered displacement-response substrate-law frontier. Once the current-body-realizing substrate body, the surjective displacement projection, the coercive positive substrate-response operator, the exact-shell discipline, and benchmark faithfulness are fixed, the \(K\)-orthogonal minimal-lift splitting is no longer separately primitive. It is uniquely forced by the canonical weighted left inverse \eqref{eq:canonical}.

That result is narrower than a completed first-principles derivation of GR from substrate microphysics, but it is stronger than another same-law deeper relay. The full displacement-response substrate law now collapses benchmark-facingly to the smaller current-body-realizing projection-operator core \(\ProjectionOperatorCore\). The same current observer-relevant pullback body-response core, pullback body-operator core, pullback-operator core, pullback-quadratic core, section-centered residual-normal-form realization, continuity-Hessian compressed core, and benchmark one-metric observer package therefore follow unchanged.

The scope remains honest. The one operative metric is still exactly \(\CompletedMetric\), the ordinary matter route is unchanged, there is no second operative metric, the low-energy matter law is not enlarged, and no stronger promotion language is licensed. The next open burden is therefore no longer the full displacement-response substrate law with a separately primitive minimal-lift family. It is to derive or force the smaller projection-operator core itself, or some nontrivial constituent or compatibility relation inside it, from still more primitive explicit substrate structure.

\end{document}
